% Source: sources/returns/S6_04_MAPS_MANIFOLDS/sections/manifolds.tex:1314–1459
\paragraph{Theorem D.2.7 (Proved: all orders $3k$ for Wang's primitive $I_2$ filling).}
Use the rational elliptic surface
\[
 y^2=x^3-3(1+t)x+(2+3t+2t^2),
 \qquad P=(1,\sqrt2\,t),\qquad Q=6P.
\]
Its fibres are $IV^*,I_2,I_1,I_1$, and the degree-one
linearization vanishes at the second $I_1$ point. Use Wang's
primitive one-component rank-one Mumford model at $I_2$, the
rank-one model at the other $I_1$, and the hexagonal rank-two
model at the zero of the linearization. In the homology marking
dual to $(\psi,u,w,\delta)$ put
\[
 \ell=e_0-2e_1-4e_2,\qquad
 m=3k,\qquad v=n\ell+r e_3,
 \qquad k>0,\quad \gcd(n,3k)=1.
\]
The order-$m$ logarithmic quotient at $IV^*$ is smooth for every
integer $r$, and the resulting threefold satisfies
\[
 \begin{gathered}
 \pi_1(X)=\mathbb Z/|n|,\qquad H_0=H_6=\mathbb Z,
 \qquad H_5=0,\\
 H_1=H_4=\mathbb Z/|n|,\qquad
 H_2=H_3=\mathbb Z/(2|n|).
 \end{gathered}
\]
In particular every $k>0$, with $n=\pm1$, gives a simply
connected complex rational homology six-sphere with two-torsion.
The case $(k,n,r)=(1,1,0)$ independently reproduces Wang's
homology theorem. The other orders are included by the proof,
without a finite parameter bound.
\begin{proof}
The source is \cite{manWang26}: Propositions 2.1 and
2.7--2.10 verify the surface and local models, and
Sections 4--5 fix their markings. In particular the $I_2$ period
exponent is $(2,1)$, not $(2,0)$. The unimodular change
$(x,y)=(v,v^2/z)$ makes it $(1,0)$ and gives a single
elliptic ruled normalization. This is the geometric input
missing from the zero-weight ``original $I_2$'' API run in D.2.2.
It is a different compactification, so the two computations
must not be equated.

Pull the smooth cubic good-reduction family back by degree $k$.
The order-three linear monodromy fixes $\ell,e_3$.
Combine base rotation of order $3k$ with translation by
$(n\ell+r e_3)/(3k)$, and use the section-preserving punctured
comparison and logarithmic clutching from D.2.6. The map $\psi$
changes by $n/(3k)$, so the action is free. This supplies the
analytic construction, properness and the marked relation
$g^{3k}=v$ for all the asserted data.

For the integral computation, use the fibre order
$(IV^*,I_2,I_1,I_1)$ and cohomological matrices
\[
 \begin{aligned}
 T_0&=\begin{pmatrix}1&0&-6&-2\\0&-1&-1&1\\0&1&0&-1\\0&0&0&1\end{pmatrix},
 &T_1&=\begin{pmatrix}1&0&6&3\\0&1&2&1\\0&0&1&0\\0&0&0&1\end{pmatrix},\\
 T_2&=\begin{pmatrix}1&0&0&0\\0&2&1&0\\0&-1&0&0\\0&0&0&1\end{pmatrix},
 &T_3&=\begin{pmatrix}1&0&0&-1\\0&1&1&0\\0&0&1&0\\0&0&0&1\end{pmatrix}.
 \end{aligned}
\]
Set $A_i=T_i^{-t}$. Their product is the identity and
$A_0^3=I$. The first two primitive vanishing directions are
$2e_2+e_3$ and $e_1+e_2$; the last filling kills $e_2,e_3$.
These generate $\langle e_1,e_2,e_3\rangle$. As in D.2.6,
the extending sections kill all three Mumford meridians and
the sphere relation kills the fourth. The remaining finite
relation is $ne_0=0$, which proves the actual fundamental-group
assertion. Before the finite filling, the same calculation
gives $\pi_1(X_D)=\mathbb Z\langle e_0\rangle$ for the
inverse image $X_D$ of a disc containing the three Mumford points.

Use $D$ and the peripheral maps $j_i$ from D.2.6. The three
Mumford invariant lattices have bases
\[
 (-e_0+3e_1,e_2,e_3),\qquad
 (e_0,e_1+e_2,e_3),\qquad(e_2,e_3).
\]
Let $K$ consist of their eight peripheral columns, followed by
$j_0(\ell)$. The submatrix on rows
\[
 (1,2,4,5,7,8,9,10,12)
\]
has determinant $-1$. Since $DK=0$ and $\operatorname{rank}D=3$, these columns
are a basis of the entire integral kernel. Also
$j_0(e_3)=-K_3-K_6-K_8$.
For the exterior matrix $[\bigwedge^2 A_i-I]$, rows
$(1,2,3,5,6)$ and columns $(1,2,4,6,7)$ have determinant $1$.
Its primitive cokernel is $\mathbb Z$, detected by
\[
 \xi=u\wedge w+6\psi\wedge\delta.
\]
The normal-circle argument of D.2.6 applies to these actual
primitive Mumford models and accounts for all eight columns.
Consequently
\[
 H_2(X_D)=Q=\mathbb Z\langle s,t\rangle,
 \quad s=[e_1\wedge e_2],\quad t=[T_{\ell,0}].
\]
Here the degree-one Mayer--Vietoris map has free domain of rank
$11$, free target of rank $12$, and cokernel $\mathbb Z$;
its kernel is zero. The stated generators therefore also
describe $H_2(X_D)$, rather than only a quotient of it.
The flat form $\xi$ has values $(1,0)$ on these generators.

In the finite linear coinvariants,
$e_1=-2e_2$ and $e_3=-3e_2$. Hence its reduced fibre has
\[
 H_1(F)=
 \mathbb Z\langle e_0,e_2,g\rangle/
 \langle mg-ne_0+3r e_2\rangle.
\]
This is free, with basis $e_2,z$, where
$z=\alpha e_0+\beta g$ and $m\alpha+n\beta=1$.
The exact norm and contraction identities are
\[
 \mathcal N_m e_2=-k e_3,\qquad
 \mathcal N_m e_0=m\ell,\qquad
 \xi(e_2,v)=2n,\qquad \xi(e_0,v)=6r.
\]
The first norm dies in $Q$. Thus the normal tori of
$e_2,e_0,g$ have coordinates $(-2n,0)$, $(-6r,m)$,
$(0,n)$. In the free basis $e_2,z$ the remaining matrix is
\[
 B_{k,n,r}=\begin{pmatrix}-2n&-6r\alpha\\0&1\end{pmatrix}.
\]
Its Smith form is $(1,2|n|)$.

All four boundary $H_1$ groups are free of total rank $14$;
the target has the same rank and cokernel $\mathbb Z/|n|$.
The degree-one kernel is zero. Each central fibre is irreducible
and a multiple is principal, so the boundary-to-filling map
in degree two is onto, by D.2.5. Therefore
$H_2(X)=\operatorname{coker}B_{k,n,r}$.
Euler localization gives $\chi(X)=2$; rational duality gives
$b_3=0$. Integral duality and the universal coefficient theorem
then give exactly the groups in the statement.

The program \texttt{code/manifolds\_wang.py} executes both
primitive-minor certificates, the norm identities and 336 full
$10\times10$ Smith forms with $1\leq k\leq12$,
$-15\leq n\leq15$, $\gcd(n,3k)=1$, and $r=0,3$.
Its input matrices and outputs are retained. The symbolic
reduction proves the formula for arbitrary displayed parameters.
\end{proof}
